\documentclass{article}

% NeurIPS 2026 Workshop: Who Verifies the Agents? (double-blind, non-archival)
\usepackage[dblblindworkshop]{neurips_2026}
\workshoptitle{Who Verifies the Agents? Toward Reliable Agent Development}

\usepackage[utf8]{inputenc}
\usepackage[T1]{fontenc}
\usepackage{hyperref}
\usepackage{url}
\usepackage{booktabs}
\usepackage{amsfonts}
\usepackage{amsmath}
\usepackage{amssymb}
\usepackage{microtype}
\usepackage{xcolor}
\usepackage{graphicx}
\usepackage{enumitem}

\newcommand{\D}{\ensuremath{\Delta}}

\title{Grounding, Not Guardrails:\\
What Actually Makes Generated Game Assets Legible}

\author{Anonymous}

\begin{document}

\maketitle

\begin{abstract}
Theme-swapping game assets with image models silently destroys gameplay legibility: under naive
prompting, validated vision--language judges recognise the gameplay role of only 42--50\% of
generated sprites, and \emph{zero percent} of damage-stage obstacles. We present a generation
pipeline grounded in a gameplay-role ontology---assets, families, and progression stages---verified
against the game engine's own tile registry, so the specification the generator follows and the
answer key the evaluator uses are the same executable object. We verify the method with
infrastructure designed to distrust itself: judges must pass a five-check validity battery before
they are allowed to score, every measurement is repeated three times, and each known evaluation
failure mode is guarded by a runnable assertion. Across 4 themes $\times$ 2 judges $\times$ 3
categories, ontology grounding never hurts (\D~$\geq 0$ in 24 of 24 cells, zero negative) and helps
most exactly where gameplay semantics cannot be inferred from pixels: obstacles $+0.417$ and
power-ups $+0.403$ versus $+0.217$ for appearance-defined match elements, with damage-stage
obstacles rising from $0.00$ to $1.00$ under both judges independently. Ablation attributes the
legibility effect to grounding alone: visual references and a critic loop contribute review-friction
control, not legibility. Grounding is necessary but not sufficient---one theme defeats it
entirely---which localises the open problem rather than hiding it.
\end{abstract}

\section{Introduction}

Generative pipelines for game assets stack mechanisms. A typical system prompts with a role
description, conditions on visual references for style coherence, runs a critic loop for quality
control, and gates output with deterministic post-processing. Such systems are evaluated end to end,
which makes them unfalsifiable in a specific way: when the full stack works, we do not learn
\emph{which} mechanism made it work, and the field accumulates mechanisms rather than
understanding.

This matters more for game assets than for images in general, because a game asset has a job. A
crate at damage stage 2 of 4 must read as \emph{that} crate, not merely as a crate; a
horizontal-clearing power-up must be distinguishable from a vertical one at a glance. These
properties are not stylistic. They are gameplay semantics, and a player who cannot read them cannot
play. We call this property \textbf{legibility}, and distinguish it from the visual consistency that
generation systems usually optimise.

Legibility is hard to measure honestly, because the obvious evaluator---ask a vision--language
model---is an opinion, not a measurement. This paper's setting removes that problem. Our test bed is
a Match-3 game whose engine contains a tile registry: executable code that defines, for every
asset, its gameplay category and its damage-stage hit points. The engine is not our opinion about
what the assets mean; it is what the game actually does. We verify our role ontology against it
before using the ontology to score anything.

With a real answer key available, we ablate the pipeline one flag at a time and find that the
legibility effect is not distributed across mechanisms. It comes from grounding generation in the
gameplay-role ontology, and it concentrates precisely on the categories whose identity is a gameplay
function rather than an appearance. The mechanisms layered on top---dual visual references, a critic
loop---do not produce a legibility gain that survives swapping the judge; what they measurably
change is review friction.

\paragraph{Related work.} Consistent multi-image generation optimises visual coherence across a
set---shared identity, style, or layout \citep{jia2025t2is,consistyle2025}---and sprite-oriented work
extends this to game-shaped outputs such as animation frames \citep{spritesheet2024}. Coherence is
necessary for an asset pack, but it is not the property a player reads: a perfectly style-consistent
pack can still be unplayable. Closest in spirit is grounding generation in a non-visual
specification, as in physically-grounded 3D asset synthesis \citep{physx3d2025} and agentic pipelines
that emit structured, engine-facing artifacts \citep{shapecraft2025,scenecraft2024}; our
specification is a gameplay-role ontology rather than physics or geometry. Multi-agent creative
systems typically compose a planner, generator, and critic \citep{crea2025,planneractorcritic2026},
and are reported end to end; we run the same composition as a single-variable ablation and find the
critic does not move the property the planner does. For evaluation, low-resolution game-art datasets
supply semantic labels annotated by people \citep{gametilenet2024}; our labels instead come from the
game engine's own tile registry, so the answer key is executable rather than annotated.

Our contributions:

\begin{itemize}[leftmargin=*,itemsep=2pt]
  \item \textbf{An executably-grounded evaluation.} The game engine's tile registry supplies the
    answer key, and we verify our ontology against it (zero category and zero stage-hit-point
    mismatches over the checkable overlap). The authority of the labels comes from shipping code,
    not from us.
  \item \textbf{An attribution result with a mechanism.} Ontology grounding never hurts
    (\D~$\geq 0$ in 24/24 cells) and helps most where gameplay semantics cannot be inferred from
    pixels (obstacles $+0.417$, power-ups $+0.403$, versus elements $+0.217$). The same experiment
    is a negative result for the two mechanisms stacked above it.
  \item \textbf{Evaluation infrastructure that distrusts itself.} Judge validity batteries,
    mandatory repeated measurement, and a runnable assertion for each failure mode we encountered.
    These are portable to any pipeline evaluated by model judges.
\end{itemize}

\section{Method: ontology-grounded asset generation}
\label{sec:method}

\paragraph{Task.} Given a target theme (e.g.\ ``ocean'') and a locked art style, generate a
complete asset pack for a Match-3 game: match elements, power-ups, obstacles, and their damage-stage
progressions. The pack must be playable, meaning a player can identify each asset's gameplay
function on sight.

\paragraph{Three mechanisms.} Our pipeline composes three, which we later ablate independently:

\begin{enumerate}[leftmargin=*,itemsep=2pt]
  \item \textbf{Ontology planning.} Generation is conditioned on a machine-readable ontology of
    gameplay roles: each asset's category (element / power-up / obstacle / background), its family,
    and its position in a progression sequence (e.g.\ crate stage 1 of 4, and which visual state
    corresponds to which remaining hit points). The generator is told what the asset \emph{does},
    not only what it should look like.
  \item \textbf{Dual visual references.} Each generation is conditioned on two images: a family
    style anchor, and the previous stage in a progression. These target style drift across a pack
    and progression collapse within a family.
  \item \textbf{Hybrid verification.} Deterministic post-processing gates (cutout integrity,
    dimensions) plus a vision--language critic scoring four axes (style, cohesion, function,
    progression), with up to three regeneration iterations.
\end{enumerate}

\paragraph{Ablation ladder.} The conditions form a single-variable ladder: \textbf{B0} has none of
the three; \textbf{B1} adds ontology planning only; \textbf{B2} adds dual references; \textbf{B3}
adds the critic loop. B0 and B1 differ in exactly one flag, with all other flags identical---a
property enforced by an assertion in the pipeline's own self-check, since the causal claim in
Section~\ref{sec:results} depends on it.

\section{Verification: an evaluation built to be falsified}
\label{sec:eval}

\subsection{The answer key precedes the question}

The game engine defines a tile registry mapping every tile identifier to its gameplay category and,
for progressions, its hit points per stage. Our ontology is hand-written, so before using it to
score anything we cross-check it against the registry: every asset present in both must agree on
category and on stage hit points. It does, with zero mismatches over the checkable overlap. The
remaining ontology entries are visual parts and stage frames with no standalone engine tile, and a
handful of engine tiles have no sprite; we report the overlap rather than implying full coverage.

This is what makes the evaluation more than an opinion poll. When we later ask a judge to name an
asset's role, the correct answer was fixed by code that ships in the game.

\subsection{Judges must qualify before they score}

A model judge is an instrument, and an uncalibrated instrument produces numbers that look like
measurements. Before a judge is allowed to produce any reported number, it must pass five checks:
\textbf{(1)} protocol compliance---answers parse as the requested format; \textbf{(2)}
test--retest agreement on repeated identical queries; \textbf{(3)} label-permutation
invariance---shuffling the answer options does not change the answer; \textbf{(4)} a blind
control---accuracy must collapse when the image is replaced by a blank, otherwise the judge is
reciting prompt priors rather than looking; \textbf{(5)} it must beat a constant-majority guess on
the sampled label distribution.

Both judges reported here pass all five. A third judge was wired in and \emph{failed} the battery,
which is the useful part: the failure was an infrastructure fault (an account without credit
returning errors that a naive harness would have silently averaged into the results). The battery
caught it; a single accuracy number would not have.

\subsection{Repeated measurement is mandatory, not optional}

Judges are stochastic. Scoring the same images twice with the same judge moved 3 of 8 numbers, by up
to $0.250$---the size of our smallest effect. Every number in Section~\ref{sec:results} is therefore
the mean of \textbf{three independent scoring passes}, with the within-cell standard deviation
reported. Measured this way, within-cell sd is $\leq 0.096$ (mean $0.025$), an order of magnitude
below the single-pass swing. Repetition, not a larger asset count, was the effective fix.

\subsection{Four failure modes, four runnable guards}

Building this evaluation produced six conclusions we later retracted, from four distinct root
causes. We report them because each is a trap available to anyone evaluating generated assets with
model judges, and because a failure mode documented in prose cannot fail---so each is now a runnable
assertion.

\begin{table}[t]
  \centering
  \caption{Evaluation failure modes encountered, and the guard that now prevents each. Every guard
  is an executable assertion that fails loudly, not a note in a README.}
  \label{tab:failures}
  \small
  \begin{tabular}{@{}p{0.24\linewidth}p{0.38\linewidth}p{0.30\linewidth}@{}}
    \toprule
    Root cause & How it misled us & Guard \\
    \midrule
    Contaminated baseline &
    A build step had overwritten the human reference art with a generated pack, inverting a
    consistency metric's apparent direction &
    Assert the reference path is distinct from, and not content-identical to, the applied pack \\
    \addlinespace
    Degenerate label space &
    Restricting candidate labels to the tiles in each level collapsed the task to a 1-of-3 choice;
    every pack scored $\geq 0.97$ regardless of quality &
    Assert the per-level candidate count is small and the global set is several times larger, so
    the degenerate setting cannot be used silently \\
    \addlinespace
    Abstention scored as error &
    A judge's refusals to answer were counted as wrong answers, turning a protocol-compliance
    problem into an apparent art-quality problem and reversing a ranking &
    Score and report abstention rate separately from accuracy; assert the two are distinguishable \\
    \addlinespace
    Under-sampling that manufactures structure &
    Single-pass scoring hid a true-zero cell behind judge noise; three items produced a rank
    agreement that vanished at five &
    Refuse to compute a per-cell effect from fewer than three passes, or a rank correlation over
    fewer than five items \\
    \bottomrule
  \end{tabular}
\end{table}

Table~\ref{tab:failures} lists them. Two points of self-criticism belong here rather than in a
footnote. First, the contaminated baseline was pointed out by a collaborator; it was not caught by
our own checks. What we can claim is that once identified, it became an assertion that fails, so it
cannot recur silently. Second, the under-sampling guards were initially only prose notes---which is
to say, no guard at all---and were made executable only after they had already misled us twice. The
headline numbers in this paper are now regenerated from the scored data by a script that refuses to
emit any claim violating these sampling rules.

\subsection{One signal that behaves}

Among all metrics we tried, judge \textbf{abstention rate} was the most useful and the least
expected. When asked to read a full game board, one judge declined to answer up to $38.9\%$ of cells
for some packs while declining $0\%$ for others. Abstention is cheap to measure, has an unambiguous
good direction, and captures something accuracy cannot: a pack that makes an observer refuse to
commit is not shippable, yet raw accuracy scores that identically to ordinary error. We recommend
reporting it as a first-class metric wherever model judges evaluate generated content.

\section{Results}
\label{sec:results}

All numbers: 4 themes, 2 validated judges, role/category identification at gameplay resolution
($70$px), mean of 3 independent scoring passes.

\subsection{Grounding is the mechanism}

\begin{table}[t]
  \centering
  \caption{Ontology grounding (B1) versus no grounding (B0). Mean of three scoring passes;
  SE from within-cell variance. Seven of eight cells clear twice their standard error; one is
  exactly zero (Section~\ref{sec:zero}); none is negative.}
  \label{tab:main}
  \small
  \begin{tabular}{@{}llccrl@{}}
    \toprule
    Theme & Judge & B0 & B1 & \D & \D~$>2$\,SE \\
    \midrule
    Fruit     & GPT-4o & 0.500 & 1.000 & $+0.500$ & yes \\
    Fruit     & Gemini & 0.417 & 0.917 & $+0.500$ & yes \\
    Pet       & GPT-4o & 0.167 & 0.667 & $+0.500$ & yes \\
    Pet       & Gemini & 0.195 & 0.583 & $+0.388$ & yes \\
    Ocean     & GPT-4o & 0.833 & 0.917 & $+0.084$ & yes \\
    Ocean     & Gemini & 0.556 & 0.917 & $+0.361$ & yes \\
    Steampunk & GPT-4o & 0.583 & 0.583 & $+0.000$ & \textbf{no} \\
    Steampunk & Gemini & 0.417 & 0.722 & $+0.306$ & yes \\
    \bottomrule
  \end{tabular}
\end{table}

Table~\ref{tab:main} shows the effect of adding ontology grounding alone. Without it, judges
recognise the gameplay role of $42$--$50\%$ of sprites on the easiest theme and as little as $17\%$
on the hardest. The single flag that tells the generator what each asset \emph{does} raises this
substantially in seven of eight cells, and never lowers it.

\subsection{Where the gain lands, and why}

\begin{table}[t]
  \centering
  \caption{Effect of grounding by ontology category, over all 24 (theme $\times$ judge $\times$
  category) cells. Every cell is non-negative. The ordering is the paper's thesis: gain tracks how
  much of an asset's identity is gameplay function rather than appearance.}
  \label{tab:bycat}
  \small
  \begin{tabular}{@{}lrrrc@{}}
    \toprule
    Category & mean \D & min & max & cells \\
    \midrule
    Obstacle (damage stages) & $+0.417$ & $+0.000$ & $+1.000$ & 8 \\
    Power-up (clear pattern) & $+0.403$ & $+0.000$ & $+0.667$ & 8 \\
    Element (colour only)    & $+0.217$ & $+0.000$ & $+0.467$ & 8 \\
    \bottomrule
  \end{tabular}
\end{table}

Table~\ref{tab:bycat} is the mechanism. \textbf{\D~$\geq 0$ in 24 of 24 cells; none is negative.}
The gain is roughly twice as large for obstacles and power-ups as for match elements, and the
headline case is damage-stage obstacles on the Fruit theme: $0.00 \to 1.00$, under both judges
independently. Without grounding, \emph{every} damage-stage crate is misread; with it, every one is
read correctly.

This ordering is not incidental. An obstacle's identity is ``stage 2 of 4, three hit points
remaining''; a power-up's is ``clears a row''. Neither is determined by appearance---many plausible
renderings of a crate are consistent with any stage. A match element's identity, by contrast, is
its colour, which the generator gets right without being told any rules. \textbf{Grounding helps
most exactly where gameplay semantics cannot be inferred from pixels}, which is the claim this paper
set out to test.

\subsection{The zero cell, and what it buys}
\label{sec:zero}

One cell is exactly zero, and it is not noise: standard deviation $0.000$ on both conditions, with
B0 and B1 getting the \emph{identical} five assets wrong across all three passes (all of them match
elements), from packs that are verifiably different images. Of the eight zero-valued
category cells, seven are ceilings---B0 already at $1.000$, leaving no headroom---and only this one
is a true zero.

We highlight it rather than bury it, for two reasons. It confirms the mechanism from the negative
side: the one place grounding does nothing is a category whose identity grounding does not encode.
And it bounds the claim. \textbf{Grounding is necessary but not sufficient}: a sufficiently
distant theme can defeat it entirely, leaving a pack whose elements no judge can read. That is a
specific, localised open problem, which is more useful to report than an averaged success.

\subsection{What references and the critic actually do}

Extending the ablation to B2 (references) and B3 (critic) does \emph{not} extend the story. Over
three seeds and two judges, neither judge's ordering of B1/B2/B3 is monotone, and the two judges
disagree with each other (Kendall $\tau = 0.333$ between judges across the four conditions). The
gaps among B1, B2 and B3 are comparable to the per-seed standard deviation. We therefore decline to
claim a legibility benefit for either mechanism.

What the critic demonstrably changes is \textbf{review friction}: it raises the needs-review rate to
$0.11 \pm 0.10$ and the mean iteration count to $1.95 \pm 0.34$, while pairwise preference between
B2 and B3 is a tie. This is a real contribution to a production pipeline---it is quality control,
catching outliers and routing them to a human---but it is a different axis from legibility, and
conflating the two would credit the critic with an effect it does not have.

\subsection{Pack ranking is not judge-stable}

A tempting use of model judges is to rank candidate asset packs. We advise against reporting such
rankings. Asking judges to identify tiles on rendered game boards, the two judges' rankings of five
packs agree at only $\tau = 0.20$. Worse, a deterministic pixel-template matcher and a
vision--language judge rank the same packs in \emph{opposite} orders: the pack that is most
pixel-distinguishable is among the least semantically legible. Pixel distinguishability and
legibility are different properties, and a metric that conflates them will rank confidently and
wrongly.

\section{Discussion}

\paragraph{Ground first, then add guardrails.} The practical recommendation is the paper's title.
Mechanisms that inject task semantics into generation paid off here; mechanisms that police the
output afterwards did not improve the property we cared about, though they did improve
manufacturability. When a pipeline underperforms on a semantic property, the productive move is to
ask what the generator was never told, before adding another verification layer.

\paragraph{Justify each mechanism by what it measurably does.} Our critic loop looked like a quality
mechanism and is in fact a triage mechanism. Both are valuable; only one was what we assumed. An
ablation that measures the right axis for each mechanism is worth more than an end-to-end number
that lets every mechanism claim credit.

\paragraph{Treat judges as instruments.} Require validity checks before scoring, require repeated
measurement, report abstention alongside accuracy, and require two judges to agree before believing
a result. Every retraction in this work would have been prevented by one of these, and each is
cheaper than the experiment it protects.

\section{Limitations}

\textbf{No human study.} Our legibility measure is model-judged. A prepared human protocol was not
run for this submission, so ``legible'' here means legible to two validated model judges, not to
players. \textbf{Small $n$ per cell.} Twelve assets per condition--theme cell; we address noise by
repeated measurement and by requiring agreement across 24 cells rather than trusting any single
cell, but the effect magnitudes are coarser than the reported decimals suggest.
\textbf{Two judges.} A third was wired in and failed the validity battery, so all judge-agreement
claims rest on two. \textbf{Single domain.} One Match-3 game; the transferable content is the
methodology, and the domain was chosen because executable ground truth is rare.
\textbf{Board-parse breadth.} The in-context board reading was run at one cell size, so we cannot
report a legibility-versus-scale curve. \textbf{One seed for three themes.} Three of the four
themes have a single generation seed; only Fruit has three.

\section{Conclusion}

Given an answer key that comes from shipping game code rather than from our own judgement, the
legibility of generated game assets is explained by one mechanism: telling the generator what each
asset does. The gain is largest precisely where gameplay semantics cannot be read off pixels, and it
is absent where they can. The mechanisms we stacked on top improve manufacturability, not
legibility---a distinction we could only draw because the evaluation was built to falsify our own
claims, and did so six times.

% Page-limit anchor: the workshop counts 4-9 pages excluding references and
% appendices, so check_paper.py reads this label's page from workshop.aux.
\label{endofbody}

\bibliographystyle{plainnat}
\bibliography{refs}

\end{document}
